\DocumentMetadata{
  pdfversion=2.0,
  pdfstandard=ua-2,
  lang=en-US,
  tagging=on,
  tagging-setup={math/setup=mathml-SE,tabsorder=structure}
}
\documentclass[11pt,letterpaper]{article}
\usepackage[margin=0.68in,includefoot]{geometry}
\usepackage{newcomputermodern}
\usepackage{amsmath,amscd,mathtools}
\usepackage{tikz,adjustbox}
\providecommand{\square}{\mdlgwhtsquare}
\usepackage[hidelinks]{hyperref}
\hypersetup{
  pdftitle={Algebra PhD Exam, May 2008},
  pdfauthor={University of Florida Department of Mathematics},
  pdfsubject={Graduate mathematics examination},
  pdfdisplaydoctitle=true,
  bookmarksopen=true
}
\setcounter{secnumdepth}{0}
\setlength{\parindent}{0pt}
\setlength{\parskip}{0.28em}
\setlength{\emergencystretch}{3em}
\setlength{\leftmargini}{2.15em}
\setlength{\leftmarginii}{2.65em}
\setlength{\leftmarginiii}{3em}
\pagestyle{empty}
\raggedbottom
\allowdisplaybreaks
\NewTaggingSocketPlug{block/list/label}{examproblem}{%
  \tagstructbegin{tag=Lbl}\tagstructend%
  \tagstructbegin{tag=\LBody}%
  \tagstructbegin{tag=\ProblemHeadingTag,title-o={\ProblemHeadingText}}%
  \tagmcbegin{tag=\ProblemHeadingTag,actualtext={\ProblemHeadingText}}%
  #2%
  \tagmcend%
  \tagstructend%
}
\newenvironment{examproblems}[2]{%
  \def\ProblemHeadingTag{#2}%
  \AssignTaggingSocketPlug{block/list/label}{examproblem}%
  \begin{enumerate}%
  \setcounter{enumi}{#1}%
  \renewcommand{\labelenumi}{\textbf{\theenumi.}}%
  \setlength{\topsep}{0.35em}%
  \setlength{\partopsep}{0pt}%
  \setlength{\itemsep}{0.48em}%
  \setlength{\parsep}{0.18em}%
}{\end{enumerate}}
\newenvironment{examsubparts}{%
  \AssignTaggingSocketPlug{block/list/label}{default}%
  \begin{enumerate}%
  \setlength{\topsep}{0.18em}%
  \setlength{\partopsep}{0pt}%
  \setlength{\itemsep}{0.16em}%
  \setlength{\parsep}{0.1em}%
}{\end{enumerate}}
\newenvironment{examinstructions}{%
  \par\begingroup\itshape
}{%
  \par\endgroup\vspace{0.18em}
}
\makeatletter
\renewcommand\subsection{\@startsection{subsection}{2}{0pt}%
  {-1.25ex plus -.25ex minus -.1ex}%
  {0.55ex plus .1ex}%
  {\normalfont\large\bfseries}}
\renewcommand{\@maketitle}{%
  \newpage\null\vskip -1em%
  {\footnotesize\bfseries
    \href{https://www.ufl.edu/}{University of Florida}, \href{https://math.ufl.edu/}{Department of Mathematics}\par}%
  \vskip -0.5em%
  \tagstructbegin{tag=H1,title={Algebra PhD Exam, May 2008}}%
  \tagpdfparaOff%
  \tagmcbegin{tag=H1}%
    {\LARGE\bfseries\href{https://gma.math.ufl.edu/exams/algebra-phd/}{Algebra PhD Exam}, May 2008\par}%
  \tagmcend%
  \tagpdfparaOn%
  \tagstructend%
  \vskip 0.25em%
  \tagmcbegin{artifact}%
    \hrule height 1pt%
  \tagmcend%
  \vskip 1em%
}
\makeatother
\title{Algebra PhD Exam, May 2008}
\author{}
\date{}
\begin{document}
\maketitle
\thispagestyle{empty}
\begin{examinstructions}
Answer seven problems. Write your answers clearly in complete English sentences. You may quote results (within reason) as long as you state them clearly.
\end{examinstructions}
\begin{examproblems}{0}{H2}
\def\ProblemHeadingText{Problem 1.}
\item
Let \(n>2\) and let~\(\zeta\) be a primitive~\(n\)-th root of unity over \(\mathbb Q\) (the field of rational numbers). Prove that
\[
[\mathbb Q(\zeta+\zeta^{-1}):\mathbb Q]=\phi(n)/2.
\]
\def\ProblemHeadingText{Problem 2.}
\item
Calculate the Galois group of \(x^5-9x+3\) over \(\mathbb Q\), the field of rational numbers. Justify your answer carefully.
\def\ProblemHeadingText{Problem 3.}
\item
Prove that the group defined by generators~\(a\) and~\(b\) and relations \(a^2=b^3=1\), \(abab=1\) is isomorphic to \(S_3\).
\def\ProblemHeadingText{Problem 4.}
\item
Let~\(A\) be a torsion abelian group. Calculate \(A\otimes_{\mathbb Z}\mathbb Q\), where \(\mathbb Q\) denotes the additive group of the rational numbers.
\def\ProblemHeadingText{Problem 5.}
\item
Define projective module. Prove from your definition that a module is projective if and only if it is a direct summand of a free module.
\def\ProblemHeadingText{Problem 6.}
\item
State and prove Hilbert’s Nullstellensatz.
\def\ProblemHeadingText{Problem 7.}
\item
Let~\(R\) be an integral domain. Define what it means for a fractional ideal to be invertible. Prove, from your definition, that if~\(I\) is an invertible fractional ideal of~\(R\), then~\(I\) is finitely generated.
\def\ProblemHeadingText{Problem 8.}
\item
Prove the following version of the Short Five Lemma. Let~\(R\) be a ring and let the following diagram of left~\(R\)-modules
{\tagpdfsetup{math/mathml/structelem=false,math/mathml/AF=false,math/alt/use=true}
\[
\begin{CD}
0 @>>> A @>{f}>> B @>{g}>> C @>>> 0 \\
@. @V{\alpha}VV @V{\beta}VV @V{\gamma}VV @. \\
0 @>>> A' @>{f'}>> B' @>{g'}>> C' @>>> 0
\end{CD}
\]
}
 be a commutative diagram whose rows are short exact sequences. Suppose that both~\(\alpha\) and~\(\gamma\) are isomorphisms. Prove that~\(\beta\) is an isomorphism.
\def\ProblemHeadingText{Problem 9.}
\item
Let~\(R\) be a commutative ring with \(1_R\). Prove that~\(R\) is a local ring if and only if for all \(r,s\in R\) such that \(r+s=1_R\), either~\(r\) or~\(s\) is a unit.
\def\ProblemHeadingText{Problem 10.}
\item
Determine up to isomorphism all semisimple noncommutative rings with \(512=2^9\) elements.
\def\ProblemHeadingText{Problem 11.}
\item
Let~\(K\) be a field and let~\(A\) be a central simple algebra over~\(K\) such that \(\dim_K(A)=n\) is finite. Prove that \(A\otimes_K A^{\mathrm{op}}\cong \mathrm{M}_{n\times n}(K)\).
\end{examproblems}
\end{document}
